         10th International Mathematical Competition for University Students
                               Cluj-Napoca, July 2003

                                                   Day 1


1. (a) Let a1 , a2 , . . . be a sequence of real numbers such that a1 = 1 and an+1 > 32 an for all n.
Prove that the sequence
                                                    an
                                                   3 n−1
                                                    
                                                    2

has a finite limit or tends to infinity. (10 points)
   (b) Prove that for all α > 1 there exists a sequence a1 , a2 , . . . with the same properties such
that
                                                  an
                                            lim n−1 = α.
                                                    3
                                                    2

(10 points)
                              an
Solution. (a) Let bn =            .   Then an+1 > 32 an is equivalent to bn+1 > bn , thus the sequence
                            3 n−1
                              
                            2

(bn ) is strictly increasing. Each increasing sequence has a finite limit or tends to infinity.
    (b) For all α > 1 there exists a sequence 1 = b1 < b2 < . . . which converges to α. Choosing
          n−1
an = 32         bn , we obtain the required sequence (an ).

2. Let a1 , a2 . . . , a51 be non-zero elements of a field. We simultaneously replace each element with
the sum of the 50 remaining ones. In this way we get a sequence b1 . . . , b51 . If this new sequence is
a permutation of the original one, what can be the characteristic of the field? (The characteristic
of a field is p, if p is the smallest positive integer such that x + x + . . . + x = 0 for any element x
                                                                  |     {z       }
                                                                            p
of the field. If there exists no such p, the characteristic is 0.) (20 points)
Solution. Let S = a1 + a2 + · · · + a51 . Then b1 + b2 + · · · + b51 = 50S. Since b1 , b2 , · · · , b51 is a

permutation of a1 , a2 , · · · , a51 , we get 50S = S, so 49S = 0. Assume that the characteristic of the
field is not equal to 7. Then 49S = 0 implies that S = 0. Therefore bi = −ai for i = 1, 2, · · · , 51.
On the other hand, bi = aϕ(i) , where ϕ ∈ S51 . Therefore, if the characteristic is not 2, the sequence
a1 , a2 , · · · , a51 can be partitioned into pairs {ai , aϕ(i) } of additive inverses. But this is impossible,
since 51 is an odd number. It follows that the characteristic of the field is 7 or 2.
     The characteristic can be either 2 or 7. For the case of 7, x1 = . . . = x51 = 1 is a possible
choice. For the case of 2, any elements can be chosen such that S = 0, since then bi = −ai = ai .

3. Let A be an n × n real matrix such that 3A3 = A2 + A + I (I is the identity matrix). Show
that the sequence Ak converges to an idempotent matrix. (A matrix B is called idempotent if
B 2 = B.) (20 points)

Solution. The minimal polynomial of A is a divisor of 3x3 − x2 − x − 1. This polynomial has three
different roots. This implies that A is diagonalizable: A = C −1 DC where D is a diagonal matrix.
The eigenvalues of the matrices A and D are all roots of polynomial 3x3 − x2 − x − 1. One of the
three roots is 1, the remaining two roots have smaller absolute value than 1. Hence, the diagonal
elements of Dk , which are the kth powers of the eigenvalues, tend to either 0 or 1 and the limit
M = lim Dk is idempotent. Then lim Ak = C −1 M C is idempotent as well.

4. Determine the set of all pairs (a, b) of positive integers for which the set of positive integers
can be decomposed into two sets A and B such that a · A = b · B. (20 points)
Solution. Clearly a and b must be different since A and B are disjoint.


                                                        1
   Let {a, b} be a solution and consider the sets A, B such that a · A = b · B. Denoting d = (a, b)
the greatest common divisor of a and b, we have a = d·a1 , b = d·b1 , (a1 , b1 ) = 1 and a1 ·A = b1 ·B.
Thus {a1 , b1 } is a solution and it is enough to determine the solutions {a, b} with (a, b) = 1.
   If 1 ∈ A then a ∈ a · A = b · B, thus b must be a divisor of a. Similarly, if 1 ∈ B, then a is a
divisor of b. Therefore, in all solutions, one of numbers a, b is a divisor of the other one.
   Now we prove that if n ≥ 2, then (1, n) is a solution. For each positive integer k, let f (k)
be the largest non-negative integer for which nf (k) |k. Then let A = {k : f (k) is odd} and
B = {k : f (k) is even}. This is a decomposition of all positive integers such that A = n · B.

5. Let g : [0, 1] → R be a continuous function and let fn : [0, 1] → R be a sequence of functions
defined by f0 (x) = g(x) and

                                  1 x
                                   Z
                       fn+1 (x) =      fn (t)dt (x ∈ (0, 1], n = 0, 1, 2, . . .).
                                  x 0

Determine lim fn (x) for every x ∈ (0, 1]. (20 points)
           n→∞
    B. We shall prove in two different ways that limn→∞ fn (x) = g(0) for every x ∈ (0, 1]. (The
second one is more lengthy but it tells us how to calculate fn directly from g.)
    Proof I. First we prove our claim for non-decreasing g. In this case, by induction, one can
easily see that
  1. each fn is non-decrasing as well, and
  2. g(x) = f0 (x) ≥ f1 (x) ≥ f2 (x) ≥ . . . ≥ g(0)            (x ∈ (0, 1]).
Then (2) implies that there exists

                                     h(x) = lim fn (x)              (x ∈ (0, 1]).
                                              n→∞

Clearly h is non-decreasing and g(0) ≤ h(x) ≤ fn (x) for any x ∈ (0, 1], n = 0, 1, 2, . . .. Therefore
to show that h(x) = g(0) for any x ∈ (0, 1], it is enough to prove that h(1) cannot be greater than
g(0).
    Suppose that h(1) > g(0). Then there exists a 0 < δ < 1 such that h(1) > g(δ). Using the
definition, (2) and (1) we get
                          Z 1                Z δ              Z 1
             fn+1 (1) =         fn (t)dt ≤         g(t)dt +         fn (t)dt ≤ δg(δ) + (1 − δ)fn (1).
                           0                  0                δ

Hence
                      fn (1) − fn+1 (1) ≥ δ(fn (1) − g(δ)) ≥ δ(h(1) − g(δ)) > 0,
so fn (1) → −∞, which is a contradiction.
    Similarly, we can prove our claim for non-increasing continuous functions as well.
    Now suppose that g is an arbitrary continuous function on [0, 1]. Let

                   M (x) = sup g(t),               m(x) = inf g(t)                  (x ∈ [0, 1])
                               t∈[0,x]                        t∈[0,x]


Then on [0, 1] m is non-increasing, M is non-decreasing, both are continuous, m(x) ≤ g(x) ≤ M (x)
and M (0) = m(0) = g(0). Define the sequences of functions Mn (x) and mn (x) in the same way
as fn is defined but starting with M0 = M and m0 = m.
    Then one can easily see by induction that mn (x) ≤ fn (x) ≤ Mn (x). By the first part of the
proof, limn mn (x) = m(0) = g(0) = M (0) = limn Mn (x) for any x ∈ (0, 1]. Therefore we must
have limn fn (x) = g(0).




                                                         2
   Proof II. To make the notation clearer we shall denote the variable of fj by xj . By definition
(and Fubini theorem) we get that
                            Z xn+1 Z xn           Z xn−1 Z x2 Z x1
                        1          1          1                   1
        fn+1 (xn+1 ) =                                   ...                g(x0 )dx0 dx1 . . . dxn
                       xn+1 0     x        xn−1 0              0 x1 0
                              ZZ n 0
                          1                                     dx0 dx1 . . . dxn
                     =                                   g(x0 )
                         xn+1    0≤x0 ≤x1 ≤...≤xn ≤xn+1            x1 . . . xn
                              Z xn+1           ZZ                                     !
                          1                                             dx1 . . . dxn
                     =               g(x0 )                                             dx0 .
                         xn+1 0                   x0 ≤x1 ≤...≤xn ≤xn+1 x1 . . . xn


   Therefore with the notation
                                                 ZZ
                                                                  dx1 . . . dxn
                                   hn (a, b) =
                                                   a≤x1 ≤...≤xn ≤b x1 . . . xn

and x = xn+1 , t = x0 we have
                                                        Z x
                                                  1
                                       fn+1 (x) =                g(t)hn (t, x)dt.
                                                  x         0

    Using that hn (a, b) is the same for any permutation of x1 , . . . , xn and the fact that the integral
is 0 on any hyperplanes (xi = xj ) we get that
                                     ZZ                                      Z b          Z b
                                                        dx1 . . . dxn                        dx1 . . . dxn
                n! hn (a, b)   =                                      =             ...
                                        a≤x1 ,...,xn ≤b x1 . . . xn           a            a  x1 . . . xn
                                       Z b     !n
                                           dx
                               =                     = (log(b/a))n .
                                        a x


   Therefore                                           Z x
                                                   1                (log(x/t))n
                                    fn+1 (x) =               g(t)               dt.
                                                   x    0                n!
  Note that if g is constant then the definition gives fn = g. This implies on one hand that we
must have
                                    1 x (log(x/t))n
                                      Z
                                                      dt = 1
                                    x 0        n!
and on the other hand that, by replacing g by g − g(0), we can suppose that g(0) = 0.
    Let x ∈ (0, 1] and ε > 0 be fixed. By continuity there exists a 0 < δ < x and an M such that
|g(t)| < ε on [0, δ] and |g(t)| ≤ M on [0, 1] . Since

                                                (log(x/δ))n
                                             lim            =0
                                            n→∞      n!
there exists an n0 sucht that log(x/δ))n /n! < ε whenever n ≥ n0 . Then, for any n ≥ n0 , we have

                                 1 x         (log(x/t))n
                                   Z
                 |fn+1 (x)| ≤         |g(t)|             dt
                                 x 0              n!
                                 1 δ (log(x/t))n          1 x        (log(x/δ))n
                                   Z                        Z
                            ≤         ε              dt +     |g(t)|             dt
                                 x 0         n!           x δ             n!
                                 1 x (log(x/t))n          1 x
                                   Z                        Z
                            ≤         ε              dt +     M εdt
                                 x 0          n!          x δ
                            ≤ ε + M ε.

   Therefore limn f (x) = 0 = g(0).


                                                             3
6. Let f (z) = an z n + an−1 z n−1 + . . . + a1 z + a0 be a polynomial with real coefficients. Prove that
if all roots of f lie in the left half-plane {z ∈ C : Re z < 0} then

                                              ak ak+3 < ak+1 ak+2

holds for every k = 0, 1, . . . , n − 3. (20 points)                            Q
Solution. The polynomial f is a product of linear and quadratic factors, f (z) = i (ki z + li ) ·

         2
Q
  j (pj z + qj z + rj ), with ki , li , pj , qj , rj ∈ R. Since all roots are in the left half-plane, for each i, ki
and li are of the same sign, and for each j, pj , qj , rj are of the same sign, too. Hence, multiplying
f by −1 if necessary, the roots of f don’t change and f becomes the polynomial with all positive
coefficients.
    For the simplicity, we extend the sequence of coefficients by an+1 = an+2 = . . . = 0 and
a−1 = a−2 = . . . = 0 and prove the same statement for −1 ≤ k ≤ n − 2 by induction.
    For n ≤ 2 the statement is obvious: ak+1 and ak+2 are positive and at least one of ak−1 and
ak+3 is 0; hence, ak+1 ak+2 > ak ak+3 = 0.
    Now assume that n ≥ 3 and the statement is true for all smaller values of n. Take a divisor of
f (z) which has the form z 2 + pz + q where p and q are positive real numbers. (Such a divisor can
be obtained from a conjugate pair of roots or two real roots.) Then we can write

                 f (z) = (z 2 + pz + q)(bn−2 z n−2 + . . . + b1 z + b0 ) = (z 2 + pz + q)g(x).                  (1)

The roots polynomial g(z) are in the left half-plane, so we have bk+1 bk+2 < bk bk+3 for all −1 ≤
k ≤ n − 4. Defining bn−1 = bn = . . . = 0 and b−1 = b−2 = . . . = 0 as well, we also have
bk+1 bk+2 ≤ bk bk+3 for all integer k.
    Now we prove ak+1 ak+2 > ak ak+3 . If k = −1 or k = n − 2 then this is obvious since ak+1 ak+2
is positive and ak ak+3 = 0. Thus, assume 0 ≤ k ≤ n − 3. By an easy computation,

                                             ak+1 ak+2 − ak ak+3 =

    = (qbk+1 + pbk + bk−1 )(qbk+2 + pbk+1 + bk ) − (qbk + pbk−1 + bk−2 )(qbk+3 + pbk+2 + bk+1 ) =
                = (bk−1 bk − bk−2 bk+1 ) + p(b2k − bk−2 bk+2 ) + q(bk−1 bk+2 − bk−2 bk+3 )+
              +p2 (bk bk+1 − bk−1 bk+2 ) + q 2 (bk+1 bk+2 − bk bk+3 ) + pq(b2k+1 − bk−1 bk+3 ).
   We prove that all the six terms are non-negative and at least one is positive. Term p2 (bk bk+1 −
bk−1 bk+2 ) is positive since 0 ≤ k ≤ n − 3. Also terms bk−1 bk − bk−2 bk+1 and q 2 (bk+1 bk+2 − bk bk+3 )
are non-negative by the induction hypothesis.
   To check the sign of p(b2k − bk−2 bk+2 ) consider

           bk−1 (b2k − bk−2 bk+2 ) = bk−2 (bk bk+1 − bk−1 bk+2 ) + bk (bk−1 bk − bk−2 bk+1 ) ≥ 0.

If bk−1 > 0 we can divide by it to obtain b2k −bk−2 bk+2 ≥ 0. Otherwise, if bk−1 = 0, either bk−2 = 0
or bk+2 = 0 and thus b2k − bk−2 bk+2 = b2k ≥ 0. Therefore, p(b2k − bk−2 bk+2 ) ≥ 0 for all k. Similarly,
pq(b2k+1 − bk−1 bk+3 ) ≥ 0.
    The sign of q(bk−1 bk+2 − bk−2 bk+3 ) can be checked in a similar way. Consider

       bk+1 (bk−1 bk+2 − bk−2 bk+3 ) = bk−1 (bk+1 bk+2 − bk bk+3 ) + bk+3 (bk−1 bk − bk−2 bk+1 ) ≥ 0.

If bk+1 > 0, we can divide by it. Otherwise either bk−2 = 0 or bk+3 = 0. In all cases, we obtain
bk−1 bk+2 − bk−2 bk+3 ≥ 0.
    Now the signs of all terms are checked and the proof is complete.




                                                         4
